\begin{table}[H]
\centering
\small
\caption{Precio FOB peruano frente a precios internacionales del Banco Mundial, 2012--2025}
\label{tab:tabla3_precios_internacionales}
\begin{tabular}{lrrrrrr}
\toprule
Producto & Referencia internacional & Precio FOB medio (US\$/kg) & Referencia media & Ratio FOB / internacional & Correlación (niveles) & Correlación (variaciones) \\
\midrule
Café & Café arábica & 3.82 & 4.28 & 0.90 & 0.99 & 0.98 \\
Cacao & Cacao & 3.50 & 3.30 & 1.07 & 1.00 & 0.97 \\
Mandarinas & Naranja & 1.16 & 0.98 & 1.35 & 0.30 & 0.47 \\
Arándanos & Índice agrícola BM & 7.87 & 100.7 & -- & -0.13 & -0.34 \\
Paltas & Índice agrícola BM & 1.88 & 100.7 & -- & -0.58 & -0.20 \\
Uvas & Índice agrícola BM & 2.47 & 100.7 & -- & 0.68 & 0.19 \\
\bottomrule
\end{tabular}
\par\vspace{2pt}\parbox{\linewidth}{\footnotesize\textit{Nota.} Promedios anuales. Precios internacionales en US\$/kg; índice agrícola con base 2010 = 100. Correlación de Pearson entre series anuales (niveles) y entre sus variaciones logarítmicas. Fuente: SUNAT (cuadro G9) y Banco Mundial (Pink Sheet); elaboración propia.}
\end{table}
